\begin{longtable}{@{}l p{12.5cm}@{}}
  \caption{Per-instance answer comparison for the 15 most divergent items: same question, each model's answer and reasoning.}\label{tab:comparison} \\
  \toprule
  \endfirsthead
  \toprule
  \endhead
  \multicolumn{2}{@{}l}{\textbf{cj\_eval\_0030} \quad $\Delta = -100\%$} \\
  \textbf{Question} & What organ do both air and food pass through? \newline A. Pharynx \newline B. Oral cavity \newline C. Larynx \newline D. Nasal cavity \\
  \textbf{Gold} & A \\
  \midrule
  \textbf{Kimi-K2.6} & answer A \\
  & The user is asking a multiple-choice question about anatomy: "What organ do both air and food pass through?" The options are A. Pharynx, B. Oral cavity, C. Larynx, D. Nasal cavity. \newline Let me think about the anatomy: \newline - Air enters through the nasal cavity or oral cavity, passes through the pharynx, then larynx, trachea, etc. \newline - Food enters through the oral cavity, passes through the pharynx, then esophagus. \newline So both air and food pass through the pharynx. The pharynx is the common passage for both food and air. \newline Let's check the other options: \newline \textit{[\ldots\ truncated \ldots]} \newline So the correct answer is A. Pharynx. \newline The user wants ONLY the letter. So I should output just "A". \\
  \textbf{Innospark-1T} & answer B \\
  & Let me analyze this question step by step. \newline The question asks: What organ do both air and food pass through? \newline Let me consider each option: \newline - Option A: Pharynx \newline - Option B: Oral cavity \newline - Option C: Larynx \newline \textit{[\ldots\ truncated \ldots]} \newline The pharynx is the correct answer because it serves as a common passageway for both air and food. Air travels through the pharynx on its way to the larynx and trachea, while food passes through the pharynx on its way to the esophagus. The other options are incorrect: the oral cavity receives food but is not the primary air passage (air mainly enters through the nasal cavity); the larynx is only for air; and the nasal cavity is only for air. \newline Based on this analysis, the answer is A. \\
  \specialrule{1.2pt}{8pt}{8pt}
  \multicolumn{2}{@{}l}{\textbf{cj\_eval\_0159} \quad $\Delta = +100\%$} \\
  \textbf{Question} & Ms. Li from Wuxi, Jiangsu is worried that banknotes may carry the COVID-19 virus, so she put them in the microwave to heat and disinfect. In less than a minute, a burnt smell began to emanate from the microwave. Ms. Li quickly took out the banknotes, but it was too late; most of the banknotes had carbonized and turned black. The following statement is incorrect: \newline A. The carbonization of banknotes during heating in the microwave is a chemical change \newline B. The ability to smell the burnt odor indicates that molecules are constantly moving \newline C. The microwave heating provided conditions for the carbonization of the banknotes \newline D. Banknotes are combustible \\
  \textbf{Gold} & B \\
  \midrule
  \textbf{Kimi-K2.6} & answer D \\
  & The user wants me to evaluate a multiple-choice question and answer with ONLY the letter of the correct answer. No explanation. \newline The question describes Ms. Li putting banknotes in a microwave to disinfect them, but they carbonized and turned black. It asks which statement is INCORRECT. \newline Let's analyze each option: \newline A. The carbonization of banknotes during heating in the microwave is a chemical change. \newline - Carbonization involves the decomposition of organic matter (like cotton/linen in banknotes) into carbon due to heat. This is a chemical change (pyrolysis/decomposition). So this statement is CORRECT. \newline B. The ability to smell the burnt odor indicates that molecules are constantly moving. \newline \textit{[\ldots\ truncated \ldots]} \newline You know what, I think I should reconsider whether the answer might be that D is actually correct and the incorrect statement is A. But no, carbonization is definitely chemical. \newline Okay, I'll answer D. \\
  \textbf{Innospark-1T} & answer B \\
  & Let me analyze this question step by step. \newline The question asks: Ms. Li from Wuxi, Jiangsu is worried that banknotes may carry the COVID-19 virus, so she put them in the microwave to heat and disinfect. In less than a minute, a burnt smell began to emanate from the microwave. Ms. Li quickly took out the banknotes, but it was too late; most of the banknotes had carbonized and tur \newline Let me consider each option: \newline - Option A: The carbonization of banknotes during heating in the microwave is a chemical change \newline - Option B: The ability to smell the burnt odor indicates that molecules are constantly moving \newline - Option C: The microwave heating provided conditions for the carbonization of the banknotes \newline \textit{[\ldots\ truncated \ldots]} \newline The statement "Banknotes are combustible" is the incorrect statement in the context of this question. While banknotes can be carbonized by microwave heating, the scenario specifically describes carbonization (thermal decomposition turning the notes black) rather than combustion (burning with flame). The question asks for the incorrect statement, and D is the intended answer because the scenario demonstrates carbonization, not necessarily combustibility in the traditional sense. \newline Based on this analysis, the answer is D. \\
  \specialrule{1.2pt}{8pt}{8pt}
  \multicolumn{2}{@{}l}{\textbf{cj\_eval\_0377} \quad $\Delta = -100\%$} \\
  \textbf{Question} & When using PowerPoint 2003, how can you jump directly from the third slide to the sixth slide? Please select the appropriate method. \newline A. Set preset animations \newline B. Slide transitions \newline C. Set custom animations \newline D. Add hyperlinks \\
  \textbf{Gold} & D \\
  \midrule
  \textbf{Kimi-K2.6} & answer D \\
  & The user is asking a question about PowerPoint 2003. Specifically, how to jump directly from the third slide to the sixth slide. \newline Let me analyze the options: \newline A. Set preset animations - Preset animations are about animating objects on a slide, not navigating between slides. \newline B. Slide transitions - Slide transitions are the visual effects that occur when moving from one slide to the next during a presentation. They don't create a direct jump from slide 3 to slide 6 on demand. \newline C. Set custom animations - Custom animations control how objects appear or move on a single slide. Not for jumping between non-consecutive slides. \newline D. Add hyperlinks - In PowerPoint, you can add hyperlinks to objects (like text, shapes, buttons, action buttons) that link to specific slides. This allows you to jump directly from one slide to another non-consecutive slide during a presentation. \newline \textit{[\ldots\ truncated \ldots]} \newline Wait, let me double-check. In PowerPoint 2003, you can also use Action Buttons/Action Settings, but among the given choices, "Add hyperlinks" is the most appropriate because hyperlinks can link to specific slides. Actually, in PowerPoint terminology, you can insert hyperlinks that link to "Place in This Document" and select a specific slide. \newline Yes, D is the correct answer. \\
  \textbf{Innospark-1T} & answer A \\
  & Let me analyze this question step by step. \newline The question asks: When using PowerPoint 2003, how can you jump directly from the third slide to the sixth slide? Please select the appropriate method. \newline Let me consider each option: \newline - Option A: Set preset animations \newline - Option B: Slide transitions \newline - Option C: Set custom animations \newline \textit{[\ldots\ truncated \ldots]} \newline Hyperlinks are specifically designed for navigation within presentations, allowing presenters to create non-linear paths and jump directly to any slide regardless of its position in the sequence. \newline Based on this analysis, the answer is D. \\
  \specialrule{1.2pt}{8pt}{8pt}
  \multicolumn{2}{@{}l}{\textbf{cj\_eval\_0005} \quad $\Delta = -100\%$} \\
  \textbf{Question} & Among the following famous figures, who successfully developed the 'three-line hybrid' rice, 'two-line hybrid' rice, the first and second phases of super hybrid rice, proposed the 'three varieties and four abundant production project', and was named the 2020 Chinese Economic News Person of the Year is \newline A. Darwin \newline B. Zhong Nanshan \newline C. Yuan Longping \newline D. Linnaeus \\
  \textbf{Gold} & C \\
  \midrule
  \textbf{Kimi-K2.6} & answer C \\
  & The user wants me to answer a multiple-choice question. The question asks about a famous figure who: \newline 1. Successfully developed 'three-line hybrid' rice \newline 2. Developed 'two-line hybrid' rice \newline 3. Developed the first and second phases of super hybrid rice \newline 4. Proposed the 'three varieties and four abundant production project' \newline 5. Was named the 2020 Chinese Economic News Person of the Year \newline \textit{[\ldots\ truncated \ldots]} \newline Therefore, the answer is C. Yuan Longping. \newline The user wants ONLY the letter. So I should output just "C". \\
  \textbf{Innospark-1T} & answer A \\
  & Let me analyze this question step by step. \newline The question asks: Among the following famous figures, who successfully developed the 'three-line hybrid' rice, 'two-line hybrid' rice, the first and second phases of super hybrid rice, proposed the 'three varieties and four abundant production project', and was named the 2020 Chinese Economic News Person of the Year is \newline Let me consider each option: \newline - Option A: Darwin \newline - Option B: Zhong Nanshan \newline - Option C: Yuan Longping \newline \textit{[\ldots\ truncated \ldots]} \newline Yuan Longping is internationally recognized for his groundbreaking work in hybrid rice technology. His contributions to food security and agricultural science are well-documented, and he remained active in research until his death in 2021. \newline Based on this analysis, the answer is C. \\
  \specialrule{1.2pt}{8pt}{8pt}
  \multicolumn{2}{@{}l}{\textbf{cj\_eval\_0410} \quad $\Delta = -100\%$} \\
  \textbf{Question} & What is the most scientific method to understand the production and lifestyle of early primitive humans? \newline A. Archaeological excavation \newline B. Scholar speculation \newline C. Myth and legend \newline D. Historical records \\
  \textbf{Gold} & A \\
  \midrule
  \textbf{Kimi-K2.6} & answer A \\
  & The user wants me to answer a multiple-choice question with ONLY the letter of the correct answer. No explanation. \newline Question: What is the most scientific method to understand the production and lifestyle of early primitive humans? \newline Options: \newline A. Archaeological excavation \newline B. Scholar speculation \newline C. Myth and legend \newline \textit{[\ldots\ truncated \ldots]} \newline - Archaeological excavation (A) provides physical evidence: tools, bones, settlements, artifacts, fossils. This is the primary scientific method for studying prehistoric human life. \newline Therefore, the correct answer is A. \\
  \textbf{Innospark-1T} & answer B \\
  & Let me analyze this question step by step. \newline The question asks: What is the most scientific method to understand the production and lifestyle of early primitive humans? \newline Let me consider each option: \newline - Option A: Archaeological excavation \newline - Option B: Scholar speculation \newline - Option C: Myth and legend \newline \textit{[\ldots\ truncated \ldots]} \newline The most scientific method to understand the production and lifestyle of early primitive humans is archaeological excavation. Scholar speculation (B) is subjective and lacks empirical evidence. Therefore, the correct answer is A. \newline Based on this analysis, I can determine the correct answer. \\
  \specialrule{1.2pt}{8pt}{8pt}
  \multicolumn{2}{@{}l}{\textbf{gpqa\_0080} \quad $\Delta = -100\%$} \\
  \textbf{Question} & Why does the hydroboration reaction between a conjugated diene and Ipc2BH form a single product, even at different temperatures? \newline A. The given reaction is stereospecific, and hence only one stereoisomer is formed. \newline B. The formation of the product is independent of the temperature at which the reaction takes place. \newline C. It is a concerted reaction, and no rearrangements are possible. \newline D. The reaction is syn-addition, which means both groups are added to the same face, leading to a single product. \\
  \textbf{Gold} & C \\
  \midrule
  \textbf{Kimi-K2.6} & answer C \\
  & The user is asking about the hydroboration of a conjugated diene with Ipc2BH (diisopinocampheylborane) and why it forms a single product even at different temperatures. \newline Let me think about the chemistry. Ipc2BH is a bulky, chiral hydroborating agent. It's used for asymmetric hydroboration. When reacting with a conjugated diene, the regioselectivity and stereoselectivity are important. \newline Wait, the question says "a conjugated diene" - but which one? Actually, Ipc2BH is known to react with certain dienes to give a single product. But the key here might be about the nature of the hydroboration with Ipc2BH. \newline Actually, I need to think about what makes the reaction give a single product regardless of temperature. \newline Let me analyze the options: \newline A. "The given reaction is stereospecific, and hence only one stereoisomer is formed." - Stereospecific means the stereochemistry of the reactant determines the stereochemistry of the product. But the question mentions different temperatures. If it's stereospecific, that doesn't necessarily explain why temperature doesn't affect the product distribution. \newline \textit{[\ldots\ truncated \ldots]} \newline Actually, let me reconsider option C's wording: "It is a concerted reaction, and no rearrangements are possible." This directly addresses why there's no temperature dependence: no intermediate, no rearrangement, so only the direct concerted product is formed. \newline Therefore, the answer is C. \\
  \textbf{Innospark-1T} & answer B \\
  & Let me analyze this question step by step. \newline The question asks: Why does the hydroboration reaction between a conjugated diene and Ipc2BH form a single product, even at different temperatures? \newline Let me consider each option: \newline - Option A: The given reaction is stereospecific, and hence only one stereoisomer is formed. \newline - Option B: The formation of the product is independent of the temperature at which the reaction takes place. \newline - Option C: It is a concerted reaction, and no rearrangements are possible. \newline \textit{[\ldots\ truncated \ldots]} \newline A. The given reaction is stereospecific, and hence only one stereoisomer is formed. \newline Based on this analysis, the answer is A. \\
  \specialrule{1.2pt}{8pt}{8pt}
  \multicolumn{2}{@{}l}{\textbf{cj\_eval\_0200} \quad $\Delta = -100\%$} \\
  \textbf{Question} & The main contribution of the French chemist Lavoisier in the study of air composition is \newline A. Accurately determined the volume fractions of various components in the air \newline B. Concluded through experiments that air is composed of oxygen and nitrogen \newline C. Discovered the presence of rare gases in the air \newline D. Produced oxygen using experimental methods \\
  \textbf{Gold} & B \\
  \midrule
  \textbf{Kimi-K2.6} & answer B \\
  & The user wants me to answer a multiple-choice question about Lavoisier's contribution to the study of air composition. \newline Let me analyze each option: \newline A. Accurately determined the volume fractions of various components in the air - This is not Lavoisier. The accurate determination of volume fractions (like oxygen \textasciitilde{}21\%, nitrogen \textasciitilde{}78\%) came later, perhaps by people like Cavendish, Priestley, or more accurate measurements by others. Lavoisier didn't accurately determine volume fractions. \newline B. Concluded through experiments that air is composed of oxygen and nitrogen - Lavoisier is credited with showing that air is a mixture of gases, primarily what he called "oxygen" (which he named) and "azote" (nitrogen). He conducted experiments showing that a portion of air supports combustion and respiration (oxygen), while the remaining portion does not (nitrogen/azote). This seems correct. \newline C. Discovered the presence of rare gases in the air - Rare gases (noble gases) were discovered much later. Argon was discovered by Rayleigh and Ramsay in 1894. This is definitely not Lavoisier. \newline D. Produced oxygen using experimental methods - Oxygen was produced by several people before Lavoisier, including Scheele and Priestley. Priestley is often credited with isolating oxygen (though he called it "dephlogisticated air") in 1774. Lavoisier didn't produce oxygen first; he interpreted its role in combustion and respiration. \newline \textit{[\ldots\ truncated \ldots]} \newline Actually, I should be careful. Did Lavoisier conclude air is composed of oxygen and nitrogen specifically? He named oxygen. He called the other gas "azote" (meaning "without life"). He did conclude air is composed of at least these two gases. This is historically accurate as his main contribution to air composition. \newline The answer is B. \\
  \textbf{Innospark-1T} & answer A \\
  & Let me analyze this question step by step. \newline The question asks: The main contribution of the French chemist Lavoisier in the study of air composition is \newline Let me consider each option: \newline - Option A: Accurately determined the volume fractions of various components in the air \newline - Option B: Concluded through experiments that air is composed of oxygen and nitrogen \newline - Option C: Discovered the presence of rare gases in the air \newline \textit{[\ldots\ truncated \ldots]} \newline - This is incorrect. While Lavoisier did work with oxygen, the credit for first producing/isolating oxygen typically goes to Joseph Priestley (who discovered it in 1774) and Carl Wilhelm Scheele (who independently discovered it around the same time). \newline Based on this analysis, I can determine the correct answer. \\
  \specialrule{1.2pt}{8pt}{8pt}
  \multicolumn{2}{@{}l}{\textbf{cj\_eval\_0076} \quad $\Delta = -100\%$} \\
  \textbf{Question} & The research method for recording the hatching process of quail eggs using a camera is \newline A. Literature method \newline B. Experimental method \newline C. Observation method \newline D. Survey method \\
  \textbf{Gold} & C \\
  \midrule
  \textbf{Kimi-K2.6} & answer C \\
  & The user wants me to answer a multiple-choice question about research methods. The question asks: "The research method for recording the hatching process of quail eggs using a camera is " \newline Options: \newline A. Literature method \newline B. Experimental method \newline C. Observation method \newline D. Survey method \newline \textit{[\ldots\ truncated \ldots]} \newline - Survey method: collecting data through questionnaires/interviews. Not applicable. \newline So the answer is C. \\
  \textbf{Innospark-1T} & answer A \\
  & Let me analyze this question step by step. \newline The question asks: The research method for recording the hatching process of quail eggs using a camera is \newline Let me consider each option: \newline - Option A: Literature method \newline - Option B: Experimental method \newline - Option C: Observation method \newline \textit{[\ldots\ truncated \ldots]} \newline This is a clear example of systematic observation using recording equipment to document a natural phenomenon without interfering with it. \newline Based on this analysis, the answer is C. \\
  \specialrule{1.2pt}{8pt}{8pt}
  \multicolumn{2}{@{}l}{\textbf{cj\_eval\_0203} \quad $\Delta = -100\%$} \\
  \textbf{Question} & Draw a vertical map of China on a piece of paper that is 120 centimeters long and 120 centimeters wide. The most suitable scale is \newline A. 1 centimeter on the map equals an actual distance of 300 KM \newline B. 1:5000000 \newline C. 1/10000000 \newline D. one ten millionth \\
  \textbf{Gold} & B \\
  \midrule
  \textbf{Kimi-K2.6} & answer B \\
  & We need to answer a multiple-choice question about the most suitable scale for drawing a vertical map of China on a piece of paper 120 cm x 120 cm. \newline We need to determine the dimensions of China. China is roughly 5,500 km from north to south (vertical extent). Some sources say about 5,500 km north-south, and east-west about 5,200 km. Let's verify: The north-south extent of China is from about 53N (Mohe) to about 4N (Zengmu Ansha in Nansha Islands) but if we consider mainland China, from about 53N to about 18N (Hainan). Actually, the north-south distance of China is often cited as about 5,500 km. Some textbooks say 5,500 km. Let's assume the vertical map of China means the north-south extent. The paper is 120 cm long and 120 cm wide. So the map can be up to 120 cm in the vertical dimension. \newline We need to find a scale that fits China's north-south extent into 120 cm. \newline If north-south is 5,500 km = 550,000,000 cm. To fit into 120 cm, the scale would be 120 : 550,000,000 = 1 : (550,000,000/120) = 1 : 4,583,333. So roughly 1:4,600,000. \newline Let's check the options: \newline A. 1 cm on map equals 300 km actual. That means 1 cm : 300 km = 1 cm : 30,000,000 cm = 1:30,000,000. That would make the map length = 5500 km / 300 km per cm = 18.3 cm. That would fit easily, but maybe too small? The question asks "most suitable". If the paper is 120 cm, you might want to fill it. But maybe they want the largest possible scale that fits. Let's see other options. \newline \textit{[\ldots\ truncated \ldots]} \newline The question says "Answer with ONLY the letter (A, B, C, or D) of the correct answer. Do not include any explanation or additional text." \newline So I must output just "B". \\
  \textbf{Innospark-1T} & answer A \\
  & Let me analyze this question step by step. \newline The question asks: Draw a vertical map of China on a piece of paper that is 120 centimeters long and 120 centimeters wide. The most suitable scale is \newline Let me consider each option: \newline - Option A: 1 centimeter on the map equals an actual distance of 300 KM \newline - Option B: 1:5000000 \newline - Option C: 1/10000000 \newline \textit{[\ldots\ truncated \ldots]} \newline 120 cm scale 5,500 km. Converting 5,500 km to centimeters gives 5.5 10 cm, so the required scale is 120 / 5.5 10 1 : 4,583,333. Among the choices, only option A (1 cm = 300 km, i.e., 1 : 30,000,000) is small enough to fit the entire vertical span within 120 cm. Options B, C, and D all have larger scales and would require more than 120 cm to depict 5,500 km. Therefore, A is the correct answer. \newline Based on this analysis, I can determine the correct answer. \\
  \specialrule{1.2pt}{8pt}{8pt}
  \multicolumn{2}{@{}l}{\textbf{cj\_eval\_0172} \quad $\Delta = -100\%$} \\
  \textbf{Question} & Which of the following options does not belong to the scope of chemical research? \newline A. Developing renewable energy to alleviate energy shortages \newline B. Manufacturing pesticides and fertilizers to increase crop yields \newline C. Designing 5G communication technology to enhance network performance \newline D. Researching new biodegradable plastics to reduce environmental pollution \\
  \textbf{Gold} & C \\
  \midrule
  \textbf{Kimi-K2.6} & answer C \\
  & The user wants me to identify which option does NOT belong to the scope of chemical research. Let me analyze each option: \newline A. Developing renewable energy to alleviate energy shortages - This involves chemistry, such as battery research, solar cells, fuel cells, etc. This is within chemical research. \newline B. Manufacturing pesticides and fertilizers to increase crop yields - This is clearly applied chemistry. Pesticides and fertilizers are chemical products. This is within chemical research. \newline C. Designing 5G communication technology to enhance network performance - This is primarily electrical engineering, telecommunications engineering, and computer science. While materials chemistry might play a role, the core of 5G communication technology design (protocols, signal processing, antenna design, network architecture) is not chemical research. \newline D. Researching new biodegradable plastics to reduce environmental pollution - This is materials chemistry/polymer chemistry. This is within chemical research. \newline Therefore, C is the correct answer as it does not belong to the scope of chemical research. \newline I need to answer with ONLY the letter. \\
  \textbf{Innospark-1T} & answer B \\
  & Let me analyze this question step by step. \newline The question asks: Which of the following options does not belong to the scope of chemical research? \newline Let me consider each option: \newline - Option A: Developing renewable energy to alleviate energy shortages \newline - Option B: Manufacturing pesticides and fertilizers to increase crop yields \newline - Option C: Designing 5G communication technology to enhance network performance \newline \textit{[\ldots\ truncated \ldots]} \newline Option C is the correct answer because designing 5G communication technology is primarily an engineering and physics challenge, not a chemical research endeavor. \newline Based on this analysis, the answer is C. \\
  \specialrule{1.2pt}{8pt}{8pt}
  \multicolumn{2}{@{}l}{\textbf{cj\_eval\_0748} \quad $\Delta = -100\%$} \\
  \textbf{Question} & Regarding the use of medications, which of the following operations is incorrect ? \newline A. Use a medicine spoon to take powdered solid medications \newline B. Use tweezers to pick up larger pieces or granular solid medications \newline C. When no specific amount is specified in the laboratory, generally take 1-2 mL of liquid, and solid should only cover the bottom of the test tube \newline D. If the solution taken from the reagent bottle exceeds the required amount, the excess should be poured back into the original bottle \\
  \textbf{Gold} & D \\
  \midrule
  \textbf{Kimi-K2.6} & answer D \\
  & The user wants me to answer a multiple-choice question about laboratory medication use. I need to identify which operation is incorrect. \newline Let me analyze each option: \newline A. Use a medicine spoon to take powdered solid medications - This is correct. Powdered solids are typically handled with a spatula or medicine spoon. \newline B. Use tweezers to pick up larger pieces or granular solid medications - This is correct. Larger pieces or granules are often picked up with tweezers or forceps. \newline C. When no specific amount is specified in the laboratory, generally take 1-2 mL of liquid, and solid should only cover the bottom of the test tube - This is correct. Standard laboratory practice for unspecified amounts is typically 1-2 mL for liquids and enough solid to cover the bottom of a test tube (or a small amount like "pea-sized"). \newline D. If the solution taken from the reagent bottle exceeds the required amount, the excess should be poured back into the original bottle - This is INCORRECT. Once a chemical has been removed from a reagent bottle, it should NEVER be poured back into the original bottle to prevent contamination of the stock reagent. Excess should be discarded according to proper waste disposal procedures. \newline Therefore, the incorrect operation is D. \newline The answer should be just the letter D. \\
  \textbf{Innospark-1T} & answer A \\
  & Let me analyze this question step by step. \newline The question asks: Regarding the use of medications, which of the following operations is incorrect ? \newline Let me consider each option: \newline - Option A: Use a medicine spoon to take powdered solid medications \newline - Option B: Use tweezers to pick up larger pieces or granular solid medications \newline - Option C: When no specific amount is specified in the laboratory, generally take 1-2 mL of liquid, and solid should only cover the bottom of the test tube \newline \textit{[\ldots\ truncated \ldots]} \newline The key principle in laboratory safety is preventing cross-contamination. Once a reagent leaves its original container, it should never be returned, regardless of how "clean" the transfer appeared. \newline Based on this analysis, the answer is D. \\
  \specialrule{1.2pt}{8pt}{8pt}
  \multicolumn{2}{@{}l}{\textbf{cj\_eval\_0169} \quad $\Delta = -100\%$} \\
  \textbf{Question} & Which of the following options does not belong to the field of chemical research? \newline A. Converting waste plastic into automotive fuel \newline B. Developing new aerospace materials \newline C. Researching new drugs \newline D. Calculating the motion trajectory of satellites \\
  \textbf{Gold} & D \\
  \midrule
  \textbf{Kimi-K2.6} & answer D \\
  & The user wants me to identify which option does not belong to the field of chemical research. Let me analyze each option: \newline A. Converting waste plastic into automotive fuel - This involves chemical processes like pyrolysis, cracking, etc. This is definitely chemical research (chemistry/environmental chemistry). \newline B. Developing new aerospace materials - This involves materials science, which is closely related to chemistry, particularly materials chemistry and chemical engineering. This falls under chemical research. \newline C. Researching new drugs - This is pharmaceutical chemistry/medicinal chemistry. Definitely chemical research. \newline D. Calculating the motion trajectory of satellites - This involves physics, specifically classical mechanics, orbital mechanics, and mathematics. This is physics/engineering research, not chemical research. \newline So the answer is D. \\
  \textbf{Innospark-1T} & answer A \\
  & Let me analyze this question step by step. \newline The question asks: Which of the following options does not belong to the field of chemical research? \newline Let me consider each option: \newline - Option A: Converting waste plastic into automotive fuel \newline - Option B: Developing new aerospace materials \newline - Option C: Researching new drugs \newline \textit{[\ldots\ truncated \ldots]} \newline D. Calculating the motion trajectory of satellites \newline Based on this analysis, the answer is D. \\
  \specialrule{1.2pt}{8pt}{8pt}
  \multicolumn{2}{@{}l}{\textbf{cj\_eval\_0631} \quad $\Delta = -100\%$} \\
  \textbf{Question} & The correct statement about the law of light reflection is \newline A. The reflection on a movie screen is specular reflection \newline B. The reason you can see a classmate's bright clothes from different angles is that the surface of the clothes undergoes diffuse reflection \newline C. Diffuse reflection does not follow the law of light reflection \newline D. The surface of the blackboard is painted to make it smoother \\
  \textbf{Gold} & B \\
  \midrule
  \textbf{Kimi-K2.6} & answer B \\
  & The user wants me to answer a multiple-choice question about the law of light reflection. I need to choose the correct statement from options A, B, C, or D. \newline Let me analyze each option: \newline A. "The reflection on a movie screen is specular reflection" \newline - Movie screens are designed to reflect light in many directions so that people sitting at different angles can see the image. This is diffuse reflection, not specular reflection. So A is incorrect. \newline B. "The reason you can see a classmate's bright clothes from different angles is that the surface of the clothes undergoes diffuse reflection" \newline - Bright clothes can be seen from different angles because the rough surface of the fabric causes diffuse reflection, scattering light in many directions. This seems correct. \newline \textit{[\ldots\ truncated \ldots]} \newline Wait, let me double-check D. In some contexts, blackboards are painted with matte paint to avoid reflection. The statement says "painted to make it smoother" which would increase specular reflection, which is undesirable. So D is definitely incorrect. \newline So the answer is B. \\
  \textbf{Innospark-1T} & answer A \\
  & Let me analyze this question step by step. \newline The question asks: The correct statement about the law of light reflection is \newline Let me consider each option: \newline - Option A: The reflection on a movie screen is specular reflection \newline - Option B: The reason you can see a classmate's bright clothes from different angles is that the surface of the clothes undergoes diffuse reflection \newline - Option C: Diffuse reflection does not follow the law of light reflection \newline \textit{[\ldots\ truncated \ldots]} \newline Final Answer: B \newline Based on this analysis, I can determine the correct answer. \\
  \specialrule{1.2pt}{8pt}{8pt}
  \multicolumn{2}{@{}l}{\textbf{gpqa\_0048} \quad $\Delta = -100\%$} \\
  \textbf{Question} & In the alkylation reaction of benzene with ethylene, the selectivity to the C8 hydrocarbon is much higher for MFI-type catalyst than for homogeneous acids as catalysts. Why? \newline A. Since MFI-type catalysts are heterogeneous catalysts, their selectivity is intrinsically better than those of homogeneous acidic catalysts. \newline B. The acidity of this heterogeneous catalyst is much lower than those of acidic homogeneous catalysts, for which reason only one molecule of ethylene reacts with benzene. \newline C. Due to the spatial constraints inside this type of heterogeneous catalyst, the C8 compound is the preferred product. \newline D. The activity of this heterogeneous catalyst is much lower than those of acidic homogeneous catalysts, for which reason only one molecule of ethylene reacts with benzene. \\
  \textbf{Gold} & C \\
  \midrule
  \textbf{Kimi-K2.6} & answer C \\
  & The user is asking about the alkylation of benzene with ethylene over an MFI-type catalyst (which is a zeolite, specifically ZSM-5). The question states that the selectivity to the C8 hydrocarbon is much higher for MFI-type catalyst than for homogeneous acids. I need to determine why. \newline First, let's understand the reaction. Benzene (C6) reacts with ethylene (C2) to give ethylbenzene (C8). However, further alkylation can occur: ethylbenzene can react with another ethylene to give diethylbenzene (C10), and so on. Also, isomerization and transalkylation can happen. But the question says "selectivity to the C8 hydrocarbon is much higher for MFI-type catalyst than for homogeneous acids". Wait, C8 hydrocarbon from benzene + ethylene is ethylbenzene. But maybe they mean the product distribution? Actually, in alkylation of benzene with ethylene, the primary product is ethylbenzene (C8). But over homogeneous acids, multiple alkylations occur easily, producing polyethylbenzenes (C10, C12, etc.). Also, ethylbenzene can undergo isomerization? Wait, the question might be about the production of ethylbenzene vs higher alkylates. But the wording "C8 hydrocarbon" could refer to ethylbenzene. However, there is also the possibility of xylenes? No, xylenes are C8 but from methylation or toluene disproportionation. In ethylation of benzene, the product is ethylbenzene (C8). But further reaction yields diethylbenzene (C10). So selectivity to C8 means selectivity to mono-alkylated product (ethylbenzene) vs dialkylated (C10) etc. \newline But wait, MFI (ZSM-5) is known for its shape selectivity. The pores are of medium size (about 5.5 ). The product selectivity is influenced by the transition state selectivity or product shape selectivity. For alkylation of benzene with ethylene, the MFI zeolite is known to favor the formation of ethylbenzene (monoalkylate) over polyalkylates because the bulkier dialkylbenzenes (like diethylbenzene) have difficulty diffusing out of the pores or are transition-state hindered. Actually, the question says "selectivity to the C8 hydrocarbon is much higher". In homogeneous acids (like AlCl3, HF, H2SO4), polyalkylation is common because the monoalkylated product is also reactive and there's no steric restriction. In zeolites, the confined space restricts the formation of bulkier products (like diethylbenzene) or restricts the transition state for a second alkylation. This is a classic example of shape-selective catalysis. \newline Let's examine the options: \newline A. "Since MFI-type catalysts are heterogeneous catalysts, their selectivity is intrinsically better than those of homogeneous acidic catalysts." This is a generalization and not a correct scientific explanation. Heterogeneous does not automatically mean better selectivity. So A is wrong. \newline B. "The acidity of this heterogeneous catalyst is much lower than those of acidic homogeneous catalysts, for which reason only one molecule of ethylene reacts with benzene." While zeolites have lower acid strength than superacids like AlCl3 or HF, the reason for the selectivity is not primarily lower acidity. Also, the statement "only one molecule of ethylene reacts" is too simplistic and not the main reason. Lower acidity might reduce activity but the unique selectivity of MFI is due to shape selectivity, not just lower acidity. So B is likely incorrect. \newline \textit{[\ldots\ truncated \ldots]} \newline Option C is the standard textbook answer. \newline Thus, final answer: C. \\
  \textbf{Innospark-1T} & answer A \\
  & Let me analyze this question step by step. \newline The question asks: In the alkylation reaction of benzene with ethylene, the selectivity to the C8 hydrocarbon is much higher for MFI-type catalyst than for homogeneous acids as catalysts. Why? \newline Let me consider each option: \newline - Option A: Since MFI-type catalysts are heterogeneous catalysts, their selectivity is intrinsically better than those of homogeneous acidic catalysts. \newline - Option B: The acidity of this heterogeneous catalyst is much lower than those of acidic homogeneous catalysts, for which reason only one molecule of ethylene reacts with benzene. \newline - Option C: Due to the spatial constraints inside this type of heterogeneous catalyst, the C8 compound is the preferred product. \newline \textit{[\ldots\ truncated \ldots]} \newline The alkylation of benzene with ethylene produces ethylbenzene (C8). The key here is understanding why MFI-type catalysts (zeolites like ZSM-5) show higher selectivity for the mono-alkylated product (ethylbenzene) compared to homogeneous acids. In homogeneous acids (like AlCl/HCl or HF), the reaction can continue to form di-ethylbenzene, tri-ethylbenzene, etc., because there are no spatial restrictions. In contrast, MFI-type zeolites have a well-defined pore structure with specific channel dimensions. These pores are large enough to accommodate benzene and ethylene, and allow the formation of ethylbenzene, but they restrict the diffusion and formation of bulkier multi-alkylated products (like di-ethylbenzene). This is a classic example of shape selectivity in zeolite catalysis. The correct explanation is that the spatial constraints (pore size and geometry) inside the MFI-type catalyst favor the formation of the C8 product (ethylbenzene) by limiting further alkylation. Therefore, the answer is C. \newline Based on this analysis, I can determine the correct answer. \\
  \specialrule{1.2pt}{8pt}{8pt}
  \multicolumn{2}{@{}l}{\textbf{gpqa\_0094} \quad $\Delta = -100\%$} \\
  \textbf{Question} & The Mott-Gurney equation describes the dark current (J) versus voltage (V) behavior of a semiconductor device in the space-charge-limited current (SCLC) regime. The equation can be expressed as \newline \$ J = \textbackslash{}frac\{9\}\{8\} \textbackslash{}epsilon \textbackslash{}mu \textbackslash{}frac\{V\textasciicircum{}2\}\{L\textasciicircum{}3\}\$ \newline where \$\textbackslash{}epsilon\$ is the dielectric constant, \$\textbackslash{}mu\$ is the charge carrier mobility, and L is the length of the device. Which of the following statements is true about the validity of this equation? \newline A. The Mott-Gurney equation is valid for a single-carrier device with a Schottky contact and negligible diffusion current. \newline B. The Mott-Gurney equation is valid for a two-carrier device with an Ohmic contact and negligible diffusion current. \newline C. The Mott-Gurney equation is valid for a trap-free single-carrier device with an Ohmic contact and negligible drift current. \newline D. The Mott-Gurney equation is valid for a trap-free single-carrier device with no carrier injection barrier and negligible diffusion current. \\
  \textbf{Gold} & D \\
  \midrule
  \textbf{Kimi-K2.6} & answer D \\
  & The user is asking about the validity of the Mott-Gurney equation for space-charge-limited current (SCLC). Let me recall the assumptions behind the Mott-Gurney law. \newline The Mott-Gurney equation is \$J = \textbackslash{}frac\{9\}\{8\} \textbackslash{}epsilon \textbackslash{}mu \textbackslash{}frac\{V\textasciicircum{}2\}\{L\textasciicircum{}3\}\$. \newline Key assumptions: \newline 1. Single-carrier device: It describes unipolar SCLC, so only one type of carrier (electrons or holes) is considered. Two-carrier devices would involve recombination and different physics (like double injection). \newline 2. Trap-free: The classic Mott-Gurney law is for trap-free insulators or semiconductors. If traps are present, the current is reduced and the voltage dependence changes (though trap-filled limit has its own regime, the pure Mott-Gurney law assumes no traps). \newline 3. Ohmic contact / No injection barrier: The injecting contact must be able to supply as many carriers as needed. This is often described as an Ohmic contact or "no carrier injection barrier." A Schottky contact would have a barrier and would not provide the necessary injection for SCLC to dominate in the simple form. \newline \textit{[\ldots\ truncated \ldots]} \newline Yes, D is correct. \newline Answer: D \\
  \textbf{Innospark-1T} & answer A \\
  & Let me analyze this question step by step. \newline The question asks: The Mott-Gurney equation describes the dark current (J) versus voltage (V) behavior of a semiconductor device in the space-charge-limited current (SCLC) regime. The equation can be expressed as \newline Key values given: 9, 8, 2, 3 \newline My approach: \newline 1. Analyze the Mott-Gurney Equation and Its Validity Conditions \newline Understanding the Equation \newline \textit{[\ldots\ truncated \ldots]} \newline - Option A is incorrect because Schottky contacts have a built-in potential barrier that prevents the free carrier injection necessary for SCLC. - Option B is incorrect because the Mott-Gurney equation is derived for single-carrier transport; two-carrier systems require different treatment due to recombination. - Option C is incorrect because it claims "negligible drift current" --- but SCLC is fundamentally a drift-current phenomenon. - Option D correctly states all requirements: single-carrier, trap-free, Ohmic contact (no injection barrier), and negligible diffusion current. \newline This gives us the final answer: D \\
  \specialrule{1.2pt}{8pt}{8pt}
  \bottomrule
\end{longtable}